%!TEX root = ../main.tex
\section{Augmentation, Synthesis, and Data Evolution}
\label{sec:evolution}

Filtering reallocates a finite pool; augmentation and synthesis attempt to create
missing coverage.  In embodied learning, a generated image is useful only when it
remains coupled to an executable state transition and a valid task.  We therefore
organize generation by what is transformed and how its validity can be tested,
not by the visual generator family.

\subsection{A Hierarchy of Generative Interventions}

\textbf{Observation-level augmentation} changes lighting, texture, background,
camera noise, crop, or viewpoint while intending to preserve state and action.
It is inexpensive and often useful for visual robustness, but must respect
occlusion and camera-frame actions.  Generating a plausible new view without
updating its extrinsics or visibility creates a mislabeled action observation.
GenAug and ROSIE use pretrained image generators to alter demonstration scenes
while preserving the intended behavior
~\cite{chen2023genaug_2302_06671,yu2023scaling_2302_11550}.  RoVi-Aug extends
the intervention to robot-video augmentation, whereas RoboEngine combines
generative editing with explicit robot-data engineering
~\cite{chen2024rovi_2409_03403,yuan2025roboengine_2503_18738}.  Together they
show the scale advantage of image-space editing; they also expose its central
audit question: whether pixels outside the edited region, geometry, action
labels, and task success remain jointly invariant.

\textbf{Trajectory-level transformation} perturbs initial states, object poses,
timing, or motion and updates actions through geometry, replay, or optimization.
It preserves more action structure than video editing but depends on accurate
object frames and task invariances.  Mirroring, retiming, and spatial warping are
valid only under the co-transformation principle from
Section~\ref{sec:acquisition}.

\textbf{Segment recomposition} decomposes demonstrations into object-relative
skills and instantiates or stitches them in new scenes.  MimicGen is the canonical
example: it extracts task segments from a small number of human demonstrations,
transforms them with respect to new object configurations, and executes the
resulting sequence to obtain additional demonstrations
~\cite{mandlekar2023mimicgen_2310_17596}.  This substantially reduces manual
demonstration demand in the paper's task suite, but assumes that task segmentation,
object pose, and segment-level invariance are reliable.  Error at a stitch can be
physically small yet place the next primitive outside its valid initiation set.

\textbf{Simulation/task generation} creates assets, scenes, goals, rewards, and
expert trajectories.  Its strength is programmatic state access and scalable
failure generation; its weaknesses are generator bias, solver artifacts, and the
reality gap.  Variation should be measured in causal task factors and contact
conditions, not only texture entropy.
GenSim composes task-generation and planning components, while RoboGen uses
generative models to propose tasks, environments, and supervision
~\cite{wang2023gensim_2310_01361,wang2023robogen_2311_01455}.  Such systems can
expand nominal task count quickly, but generated tasks inherit the proposer,
asset library, reward generator, and planner's joint support; counting tasks is
therefore not evidence of behavioral coverage.

\textbf{Video and world-model generation} predicts demonstrations or counterfactual
futures conditioned on language, actions, images, or geometry.  Xiaomi-Robotics-U0
targets embodied synthesis with multi-view and geometric constraints
~\cite{li2026xiaomi_2607_11643}.  Such models can propose rare scenes and future
states at much lower marginal cost than real collection, but a photorealistic
rollout can hide inconsistent 3-D geometry, unreachable motion, or incorrect
contact.  Generated pixels do not automatically supply calibrated robot actions;
action recovery is a separate inverse problem whose uncertainty must be stored.

\subsection{The Bottleneck Is Layered Validation}

Generation quality is sometimes evaluated with visual fidelity or a
vision--language score.  These are necessary at most, never sufficient.  The
validation stack in Table~\ref{tab:validation-stack} proceeds from cheap technical
checks to expensive deployment evidence.  A sample should advance only as far as
its intended training role requires: an observation-only pretraining clip may not
need action validity, whereas a contact-rich behavior-cloning target must pass
kinematic and physical gates.

\begin{table}[t]
\centering
\caption{A validation stack for transformed or generated embodied data.  Passing
an upper row does not imply passing a lower row.}
\label{tab:validation-stack}
\small
\setlength{\tabcolsep}{4pt}
\begin{tabularx}{\columnwidth}{@{}p{1.8cm}X p{1.8cm}@{}}
\toprule
Layer & Typical question & Example signal \\
\midrule
Technical & Is the record readable and synchronized? & Schema, timestamps, missingness \\
Semantic & Does the record express the requested task? & VLM judgment, label consistency \\
Geometric & Are pose, scale, visibility, and collision plausible? & Reconstruction, collision checks \\
Kinematic & Can the target embodiment realize the motion? & IK, joint and velocity limits \\
Dynamic / contact & Are forces and interactions physically credible? & Simulation replay, contact constraints \\
Policy & Does training on it improve behavior? & Fixed-budget rollout utility \\
Deployment & Does the gain survive the target distribution? & Real success, recovery, safety \\
\bottomrule
\end{tabularx}
\end{table}


The layers have asymmetric errors.  A semantic verifier can accept a visually
correct but physically impossible trajectory; inverse kinematics can find a
joint path that violates force or collision constraints; a simulator can replay a
trajectory under inaccurate contact parameters; and a policy can exploit a
simulator artifact.  Agreement between generator and verifier is also weak when
they share training data or architecture.  High-risk samples should be checked by
heterogeneous models, analytic constraints, and real measurements where possible.

Validation should emit calibrated uncertainty and reason codes.  Hard rejection
throws away informative near misses that can train a verifier, world model, or
recovery policy.  A better routing policy sends: valid action supervision to
policy training; semantically valid but action-uncertain video to representation
learning; physically invalid samples to verifier negatives; and ambiguous
high-value cases to simulation or human review.  The relevant quantity is
\emph{verified yield per unit cost}, not raw generation throughput.

\subsection{Generative Diversity Can Be Illusory}

Counting prompts or rendered scenes overstates coverage when samples share a
generator's narrow prior.  Diversity should be measured at several levels:
visual appearance, object and layout factors, action paths, state transitions,
task strategy, failure cause, and embodiment.  Near-duplicate video may still
contain distinct contact dynamics, while visually diverse generations may induce
the same action sequence.  Joint state--action or transition-space diagnostics
are therefore more relevant to control than image embeddings alone, an intuition
also reflected in SCIZOR's joint-space deduplication
~\cite{zhang2025scizor_2505_22626}.

Generated data also changes the learner that later evaluates it.  A verifier
calibrated on an earlier policy can overvalue already mastered synthetic states or
miss novel failure modes.  Periodic controlled retraining and rollout audits are
needed to recalibrate generation priorities.  The ideal generator and verifier
form a constrained adversarial relationship: the generator targets coverage gaps
and challenging valid cases, while independent verifiers search for semantic and
physical violations rather than rewarding familiar-looking output.
